% ============================================================
% Axis1Review.tex — Artículo IEEE, Etapa 1 (preliminar)
% Build: latexmk -pdf Axis1Review.tex
% ============================================================
\documentclass[conference]{IEEEtran}

% ============================================================
% preamble.tex — configuración compartida, compatible con IEEEtran
%
% NO agregar: geometry, fancyhdr, setspace, tocloft, natbib,
% caption, subcaption, cleveref. IEEEtran los maneja internamente.
% ============================================================

% ---------- Idioma ----------
% es-tabla     -> "Tabla" en vez de "Cuadro"
% es-noquoting -> evita que babel haga activo el carácter " (rompe listings)
% es-nodecimaldot -> respeta el punto decimal si se usa siunitx
\usepackage[T1]{fontenc}
\usepackage[spanish,es-tabla,es-noquoting,es-nodecimaldot]{babel}
\usepackage{microtype}

% ---------- Encabezados de IEEEtran en español ----------
\renewcommand{\abstractname}{Resumen}
\renewcommand{\IEEEkeywordsname}{Palabras clave}
\renewcommand{\refname}{Referencias}
\renewcommand{\figurename}{Fig.}
\renewcommand{\tablename}{TABLA}

\usepackage{amsmath}
\usepackage{amssymb}
\usepackage{amsthm}
\usepackage{mathtools}
\usepackage{siunitx}
\sisetup{per-mode=symbol}

\usepackage{cite}

\usepackage{xcolor}
\definecolor{tecnavy}{RGB}{31,56,100}
\definecolor{tecsteel}{RGB}{46,85,150}
\definecolor{tecrule}{RGB}{155,184,216}
\definecolor{boxbg}{RGB}{242,246,252}
\definecolor{labelgray}{RGB}{120,120,120}
\definecolor{codebg}{RGB}{248,248,248}
\definecolor{codekw}{RGB}{0,0,180}
\definecolor{codecmt}{RGB}{0,120,60}
\definecolor{codestr}{RGB}{170,55,55}

\usepackage{graphicx}
\graphicspath{{img/}{figures/}{./}}
\usepackage{tikz}
\usepackage{pgfplots}
\pgfplotsset{compat=1.18}
\usetikzlibrary{arrows.meta,positioning,shapes.geometric,calc,fit,
                backgrounds,decorations.pathreplacing}

\usepackage{array}
\usepackage{booktabs}
\usepackage{multirow}
\usepackage{makecell}
\usepackage{tabularx}
\usepackage{threeparttable}
\usepackage{enumitem}

\newcolumntype{Y}{>{\raggedright\arraybackslash}X}
\newcolumntype{P}[1]{>{\raggedright\arraybackslash}p{#1}}
\renewcommand{\arraystretch}{1.15}
\setlist[itemize]{leftmargin=*,noitemsep,topsep=2pt}
\setlist[enumerate]{leftmargin=*,noitemsep,topsep=2pt}

\theoremstyle{definition}
\newtheorem{definicion}{Definición}
\newtheorem{ejemplo}{Ejemplo}
\theoremstyle{plain}
\newtheorem{proposicion}{Proposición}
\theoremstyle{remark}
\newtheorem*{observacion}{Observación}

\usepackage{listings}
\lstdefinestyle{code}{
  basicstyle=\ttfamily\scriptsize,
  keywordstyle=\color{codekw}\bfseries,
  commentstyle=\color{codecmt}\itshape,
  stringstyle=\color{codestr},
  backgroundcolor=\color{codebg},
  numbers=left,
  numberstyle=\tiny\color{labelgray},
  numbersep=6pt,
  frame=single,
  rulecolor=\color{tecrule},
  breaklines=true,
  postbreak=\mbox{\textcolor{tecsteel}{$\hookrightarrow$}\space},
  showstringspaces=false,
  tabsize=2,
  captionpos=b,
  columns=flexible,
  upquote=true
}
\lstset{style=code}

\usepackage[hidelinks]{hyperref}

% ---------- Entrada segura de archivos, imágenes y código ----------
\newcommand{\cajaarchivofaltante}[1]{%
  \par\medskip\noindent
  \fcolorbox{red!60!black}{red!5}{%
    \parbox{\dimexpr\linewidth-2\fboxsep-2\fboxrule}{%
      \centering\ttfamily\scriptsize ARCHIVO FALTANTE: \detokenize{#1}}}%
  \par\medskip}

\newcommand{\cajaimagenfaltante}[1]{%
  \begingroup\setlength{\fboxsep}{0pt}%
  \fcolorbox{tecrule}{boxbg}{%
    \parbox[c][2.4cm][c]{0.9\linewidth}{\centering\ttfamily\scriptsize
      IMAGEN NO ENCONTRADA\\[3pt]\detokenize{#1}}}%
  \endgroup}

\newcommand{\safeinput}[1]{%
  \IfFileExists{#1.tex}{\input{#1}}{%
    \PackageWarning{tecpaper}{Falta el archivo de sección #1.tex}%
    \cajaarchivofaltante{#1.tex}}}

\newcommand{\safeimage}[2][width=\columnwidth]{%
  \IfFileExists{#2}{\includegraphics[#1]{#2}}{%
    \IfFileExists{img/#2}{\includegraphics[#1]{img/#2}}{%
      \PackageWarning{tecpaper}{Falta la imagen #2}%
      \cajaimagenfaltante{#2}}}}

\newcommand{\code}[1]{\texttt{#1}}

% ---------- Marcador de pendiente ----------
% \pc = "por completar". Se imprime en rojo para que ningún hueco
% pase inadvertido. ANTES DE ENTREGAR, redefinirlo como vacío:
%   \newcommand{\pc}[1]{}
% y recompilar: todo lo que desaparezca nunca se escribió.
\newcommand{\pc}[1]{%
  \texorpdfstring{\textcolor{red}{\textbf{[POR COMPLETAR: #1]}}}%
                 {[POR COMPLETAR: #1]}}

% Numeración de páginas — no es el comportamiento por defecto de
% IEEEtran conference. Se activa porque este documento se entrega
% como trabajo de curso, no como sumisión a una conferencia real.
\pagestyle{plain}

\begin{document}

% ------------------------------------------------------------
% TÍTULO
% Descriptivo: técnica + dominio + dataset. Sin siglas oscuras.
% ------------------------------------------------------------
\title{\pc{Título del artículo --- técnica + dominio + dataset}\\
{\footnotesize \textit{Versión preliminar --- Etapa 1: estado del arte y diseño propuesto}}}

\author{
\IEEEauthorblockN{Adriel S. Chaves Salazar}
\IEEEauthorblockA{\textit{Escuela de Ingenier\'ia en Computaci\'on} \\
\textit{Instituto Tecnol\'ogico de Costa Rica}\\
Cartago, Costa Rica \\
2021031465}
\and
\IEEEauthorblockN{Daniel Duarte Cordero}
\IEEEauthorblockA{\textit{Escuela de Ingenier\'ia en Computaci\'on} \\
\textit{Instituto Tecnol\'ogico de Costa Rica}\\
Cartago, Costa Rica \\
{\pc{carné}}}
\and
\IEEEauthorblockN{Sebasti\'an Hern\'andez Bonilla}
\IEEEauthorblockA{\textit{Escuela de Ingenier\'ia en Computaci\'on} \\
\textit{Instituto Tecnol\'ogico de Costa Rica}\\
Cartago, Costa Rica \\
{\pc{carné}}}
}

\maketitle

% ------------------------------------------------------------
% RESUMEN PRELIMINAR
% Se marca de forma EXPLÍCITA como preliminar.
% Describe problema, motivación y enfoque. NO reporta resultados.
% ------------------------------------------------------------
\begin{abstract}
\textbf{Resumen preliminar.} Este resumen describe el problema, la
motivación y el enfoque propuesto. No incorpora resultados de ningún
tipo, dado que en esta etapa no se ha ejecutado ningún experimento.

\pc{1--2 oraciones: el problema y por qué importa}
\pc{1 oración: la brecha identificada en la literatura}
\pc{1--2 oraciones: el enfoque propuesto y el track elegido}
\pc{1 oración: el protocolo de evaluación y el baseline de literatura
que el trabajo busca igualar o superar}
\end{abstract}

% ------------------------------------------------------------
% PALABRAS CLAVE
% 4 a 6. Términos por los que alguien buscaría este trabajo.
% ------------------------------------------------------------
\begin{IEEEkeywords}
\pc{palabra clave}, \pc{palabra clave}, \pc{palabra clave},
\pc{palabra clave}, \pc{palabra clave}
\end{IEEEkeywords}

\safeinput{sec/01introduccion}
\safeinput{sec/02trabajos_relacionados}
\safeinput{sec/03baseline_literatura}
\safeinput{sec/04metodologia}
\safeinput{sec/05analisis_tecnico}
\safeinput{sec/06etica}

\bibliographystyle{IEEEtran}
\bibliography{sec/99referencias}

\end{document}